\section{Dados e recorte da análise}
\label{sec:dados}

\subsection{Fontes, unidades e preparação}

A base econômica é a MIP brasileira de 2015, disponibilizada pelo IBGE no detalhamento de 67 atividades \parencite{ibge2015mip}. A atividade produtiva é a unidade de análise, e os valores monetários são expressos em milhões de reais, a preços básicos de 2015.

A Tabela~\ref{tab:fontes} reúne as fontes, dimensões e unidades das entradas. A demanda final inclui usos domésticos e exportações de produtos nacionais. O vetor por atividade é obtido distribuindo a demanda dos 127 produtos entre suas atividades produtoras, conforme as participações da Tabela 13 do IBGE.

\begin{table}[htbp]
\centering
\caption{Fontes e objetos de entrada da análise.}
\label{tab:fontes}
\small
\begin{tabular}{@{}>{\raggedright\arraybackslash}p{0.26\textwidth}>{\raggedright\arraybackslash}p{0.19\textwidth}>{\raggedright\arraybackslash}p{0.46\textwidth}@{}}
\toprule
Objeto & Dimensão & Fonte e unidade \\
\midrule
Produção bruta ($x$) & $67\times1$ & IBGE, Tabela 01; R\$ milhões \\
Demanda final por atividade ($y$) & $67\times1$ & IBGE, Tabela 03, convertida com a Tabela 13; R\$ milhões \\
Coeficientes técnicos ($A$) & $67\times67$ & IBGE, Tabela 14; adimensional \\
Inversa de Leontief ($L$) & $67\times67$ & IBGE, Tabela 15; adimensional \\
Intensidades diretas ($\varphi^{2011}$, $\varphi^{2018}$) & $67\times1$ por ano & Sanguinet e Azzoni (2024), Tabela 2; Gg CO\textsubscript{2}/R\$ milhão \\
\bottomrule
\end{tabular}
\par\smallskip
\begin{minipage}{\textwidth}\footnotesize
Fonte: IBGE (2018) e Sanguinet e Azzoni (2024). Os anos ambientais são 2011 e 2018; a base econômica refere-se a 2015. As dimensões descrevem a organização das entradas, não observações estatisticamente independentes.
\end{minipage}
\end{table}

A Tabela 2 de \textcite{sanguinetAzzoni2024} fornece intensidades diretas de CO\textsubscript{2} por unidade de produção bruta para 2011 e 2018. A correspondência nominal com as 67 atividades da MIP é direta, sem agregação ou desagregação local. Utiliza-se um vetor por ano, sem intensidades observadas separadamente para cada estado.

Para o cenário de 2015, utiliza-se interpolação linear:
\begin{equation}
\varphi_i^{\mathrm{cen},2015}=\frac{3}{7}\varphi_i^{2011}+\frac{4}{7}\varphi_i^{2018}.
\label{eq:coeficiente-cenario}
\end{equation}
A interpolação pressupõe evolução linear da intensidade entre os dois anos e não incorpora observação física de emissões em 2015.\footnote{A fonte publica duas casas decimais; casas adicionais do cálculo não representam precisão ambiental adicional.}

\subsection{Fronteira e validade ambiental}

O recorte abrange emissões de CO\textsubscript{2} das atividades produtivas brasileiras para demanda doméstica e exportações. Exclui emissões estrangeiras na produção de importações e no processamento ou uso posterior dos produtos exportados. O destino da atribuição identifica o setor do produto final na fronteira nacional. A conta pela demanda não corresponde, portanto, à pegada mundial do consumo dos residentes.

A análise contabiliza CO\textsubscript{2}, sem agregar os demais gases de efeito estufa em CO\textsubscript{2} equivalente. A cobertura ambiental permanece incerta: a documentação da fonte apresenta descrições divergentes sobre uso da terra, e não foram recuperados os numeradores físicos nem os procedimentos completos de alocação das emissões. A correspondência setorial não certifica a cobertura nem o volume físico atribuído a cada atividade.

A compatibilidade monetária também permanece pendente. A fonte relata atualização de sua matriz inter-regional para preços de dezembro de 2018, enquanto esta aplicação combina coeficientes interpolados com valores de 2015, sem ajuste adicional. A igualdade das unidades não resolve essa diferença de preços; é necessário verificar os denominadores setoriais antes de interpretar os volumes estimados.

A base ambiental constitui, assim, um \textbf{cenário exploratório}. As verificações de balanço e das identidades contábeis asseguram consistência dos cálculos, mas não validam as emissões físicas. Conclusões ambientais definitivas dependem da reconstrução e validação dos coeficientes ou de sua substituição por uma extensão compatível com 2015.